\section{Conclusion}
\label{sec:conclusion}
\guide{About 450 words. Explain how the alliance structure shapes ecosystem
coordination (Task 1), network effects and lock-in (Task 2), platform
governance and pricing (Task 3), cybersecurity incentives (Task 4) and value
creation and capture (Task 5). Pull one finding from each task.}
\fillin{Text}

\guide{End with the single most important strategic dependency or tension, in
one clear sentence, then two or three sentences on why it matters.}
\textbf{Key strategic tension:} \fillin{One sentence}

\textbf{LT FULLFØRER HER NÅR ALLE SEKSJONENE ER FERDY}

AI was useful to us in three ways. It gave a fast first overview of a case that none of us knew well, it produced a checklist of actors and concepts to test, and it forced us to state our own reasoning clearly, because we had to explain exactly why its answer was wrong. Where the AI was right, for example on the main roles of Thales, Google and ANSSI in Task 1, it saved time.

Its failures followed a clear pattern. First, it tended to describe the alliance as if one firm owned it. In Task 1 it placed S3NS at the centre as a hub and described the dependency between Thales and Google as balanced, while the evidence shows a one-way, co-specialised dependency (Appendix A.1). In Task 3 it [fill in: e.g. listed Google as a platform side / ignored who funds the subsidy]. Second, it gave precise answers where the honest answer is ``not disclosed''. [Fill in: invented prices in Task 3, probabilities or ROI figures in Task 4, if they appeared.] Third, it used course concepts loosely. [Fill in from Task 2: e.g. economies of scale labelled as network effects.] [Fill in from Task 5: e.g. canvas blocks with no split between partners.]

These failures were not random. The AI works from what is typical for cloud companies and platforms in general, and it fills gaps with plausible content. Our case depends on what is specific: a French state that treats sovereignty as a legal requirement, a partner whose technology cannot be replaced, and costs and risks that are split unevenly between the partners. Economic judgment means deciding which theory applies, which facts are missing, and how much weight a claim can carry. The AI did not do this on its own. It could only help once we had written our own analysis first and had a standard to measure it against.

Our main lesson is that AI is a good sparring partner but a poor analyst. It can speed up the search for ideas, but every claim it made had to be checked against the textbook, the literature and the case sources before we could use it.
